\documentclass[10pt,a4paper]{amsart}
\usepackage{amsmath,amssymb,amsthm,graphicx}
% Retained preamble for the Stacks Project chapter snapshot in this directory.
% The chapter's own preamble.tex is a copy of the shared Stacks Project preamble; the block
% below is sliced MECHANICALLY, line for line, from it. Nothing was added, removed or
% redefined; the wrapper class, the xr-hyper/hyperref wiring and the externaldocument list
% are omitted because the assembler supplies the document class.
%
% lines taken from preamble.tex: 3, 5, 9, 11, 12, 156, 157, 158, ... (51 lines total)
%
% package coverage restored from the chapter's own preamble (2026-10-13): xy, xy-2cell

\usepackage[all]{xy}
\xyoption{2cell}
\UseAllTwocells

\newenvironment{reference}{\comment}{\endcomment}
\newenvironment{slogan}{\comment}{\endcomment}
\newenvironment{history}{\comment}{\endcomment}
\theoremstyle{plain}
\newtheorem{theorem}[subsection]{Theorem}
\newtheorem{proposition}[subsection]{Proposition}
\newtheorem{lemma}[subsection]{Lemma}
\theoremstyle{definition}
\newtheorem{definition}[subsection]{Definition}
\newtheorem{example}[subsection]{Example}
\newtheorem{exercise}[subsection]{Exercise}
\newtheorem{situation}[subsection]{Situation}
\theoremstyle{remark}
\newtheorem{remark}[subsection]{Remark}
\newtheorem{remarks}[subsection]{Remarks}
\numberwithin{equation}{subsection}
\def\lim{\mathop{\mathrm{lim}}\nolimits}
\def\colim{\mathop{\mathrm{colim}}\nolimits}
\def\Spec{\mathop{\mathrm{Spec}}}
\def\Hom{\mathop{\mathrm{Hom}}\nolimits}
\def\Ext{\mathop{\mathrm{Ext}}\nolimits}
\def\SheafHom{\mathop{\mathcal{H}\!\mathit{om}}\nolimits}
\def\SheafExt{\mathop{\mathcal{E}\!\mathit{xt}}\nolimits}
\def\Sch{\mathit{Sch}}
\def\Mor{\mathop{\mathrm{Mor}}\nolimits}
\def\Ob{\mathop{\mathrm{Ob}}\nolimits}
\def\Sh{\mathop{\mathit{Sh}}\nolimits}
\def\NL{\mathop{N\!L}\nolimits}
\def\CH{\mathop{\mathrm{CH}}\nolimits}
\def\proetale{{pro\text{-}\acute{e}tale}}
\def\etale{{\acute{e}tale}}
\def\QCoh{\mathit{QCoh}}
\def\Ker{\mathop{\mathrm{Ker}}}
\def\Im{\mathop{\mathrm{Im}}}
\def\Coker{\mathop{\mathrm{Coker}}}
\def\Coim{\mathop{\mathrm{Coim}}}
\DeclareMathSymbol{\boxtimes}{\mathbin}{AMSa}{"02}
\def\QCohstack{\mathcal{QC}\!\mathit{oh}}
\def\Cohstack{\mathcal{C}\!\mathit{oh}}
\def\Spacesstack{\mathcal{S}\!\mathit{paces}}
\def\Quotfunctor{\mathrm{Quot}}
\def\Hilbfunctor{\mathrm{Hilb}}
\def\Curvesstack{\mathcal{C}\!\mathit{urves}}
\def\Polarizedstack{\mathcal{P}\!\mathit{olarized}}
\def\Complexesstack{\mathcal{C}\!\mathit{omplexes}}
\def\Pic{\mathop{\mathrm{Pic}}\nolimits}
\def\Picardstack{\mathcal{P}\!\mathit{ic}}
\def\Picardfunctor{\mathrm{Pic}}
\def\Deformationcategory{\mathcal{D}\!\mathit{ef}}
\usepackage{cleveref}
\crefname{originaltheorem}{original theorem}{original theorems}
\Crefname{originaltheorem}{Original theorem}{Original theorems}
\makeatletter
\newlabel{definition-covering-SR}{{01G3}{source}{}{sourcepointer-1}{}}
\newlabel{definition-covering-SR@cref}{{[originaltheorem][0][]01G3}{[1][0][]source}{}{}{}}
\makeatother
\crefname{originaltheorem}{original theorem}{original theorems}
\Crefname{originaltheorem}{Original theorem}{Original theorems}
\makeatletter
\newlabel{example-cech}{{01G6}{source}{}{sourcepointer-2}{}}
\newlabel{example-cech@cref}{{[originaltheorem][0][]01G6}{[1][0][]source}{}{}{}}
\makeatother
\crefname{originaltheorem}{original theorem}{original theorems}
\Crefname{originaltheorem}{Original theorem}{Original theorems}
\makeatletter
\newlabel{lemma-cech-spectral-sequence}{{01GY}{source}{}{sourcepointer-3}{}}
\newlabel{lemma-cech-spectral-sequence@cref}{{[originaltheorem][0][]01GY}{[1][0][]source}{}{}{}}
\makeatother
\crefname{originaltheorem}{original theorem}{original theorems}
\Crefname{originaltheorem}{Original theorem}{Original theorems}
\makeatletter
\newlabel{lemma-covering}{{01GJ}{source}{}{sourcepointer-4}{}}
\newlabel{lemma-covering@cref}{{[originaltheorem][0][]01GJ}{[1][0][]source}{}{}{}}
\makeatother
\crefname{originaltheorem}{original theorem}{original theorems}
\Crefname{originaltheorem}{Original theorem}{Original theorems}
\makeatletter
\newlabel{lemma-h0-cech}{{01GV}{source}{}{sourcepointer-5}{}}
\newlabel{lemma-h0-cech@cref}{{[originaltheorem][0][]01GV}{[1][0][]source}{}{}{}}
\makeatother
\crefname{originaltheorem}{original theorem}{original theorems}
\Crefname{originaltheorem}{Original theorem}{Original theorems}
\makeatletter
\newlabel{sites-cohomology-lemma-kill-cohomology-class-on-covering}{{01FW}{source}{}{sourcepointer-6}{}}
\newlabel{sites-cohomology-lemma-kill-cohomology-class-on-covering@cref}{{[originaltheorem][0][]01FW}{[1][0][]source}{}{}{}}
\makeatother
\title{Sheaf cohomology agrees with hypercovering Cech cohomology}
\author{The Stacks Project Authors}
\date{stacks-project:hypercovering.tex:local-snapshot-20261012; GNU FDL 1.2 or later}
\begin{document}
\maketitle
\noindent\textit{Editorial note.} Statement span is lines 2205-2216 of sources/stacks-hypercovering/pinned-original-source.tex; proof span is lines 2218-2309, the complete author proof (90 body lines).\par\medskip
\noindent\textit{Editorial note.} Navigation stubs: 6 cited-but-not-included labels are supplied as non-load-bearing pointers with their pinned tags where the snapshot tag index has them.\par\medskip
\noindent\textit{Editorial note.} extra\_packages is empty on purpose: the retained preamble already loads xy with options and a second declaration clashes in the assembler wrapper.\par\medskip
\noindent\textit{Editorial note.} Text is retained verbatim, including the author's own typos and placeholders.\par\medskip
\noindent\textit{Editorial note.} Queue index 3 of queue\_hypercovering.json (pb=4332, stmt 2205-2216, proof 2218-2309).\par\medskip
\noindent\textit{Editorial note.} stacks-hypercovering wave-1 (hc01): fresh adjudication accept (2026-10-14); pool per standing fallback order (hypercovering -> pic -> spaces-more-cohomology) after md01 closed the moduli pool; hypercovering.tex 53 labels/34 with proof judged 26 accept / 8 reject / 14 pb=0 env-dup rows excluded; spans cite-free and reference/slogan/history-free so no patchenv needed.\par\medskip
\noindent\textit{Editorial note.} Cross-chapter same-label warning (NOT an occupancy conflict, Chief ruling 2026-10-13 evening): label 'theorem-cohomology-hypercoverings' also exists in stacks-snapshot-20261011 — a different proposition in a different chapter file.\par\medskip
\noindent\textit{Allowed standard prerequisites (editorial).}\begin{itemize}
\item Explicit prerequisite boundary. The cited inputs are prerequisites; the argument in the retained span is not replaced by any cited result.
\end{itemize}\medskip
\ifdefined\hypertarget\hypertarget{sourcepointer-1}{}\fi
\noindent\textit{Source reference.} Original theorem 01G3: sources/stacks-hypercovering/pinned-original-source.tex, source lines 510--521. Cited by the retained statement or proof but not part of the collected spans; supplied only as a navigation pointer. The argument does not depend on its content.\par\medskip
\ifdefined\hypertarget\hypertarget{sourcepointer-2}{}\fi
\noindent\textit{Source reference.} Original theorem 01G6: sources/stacks-hypercovering/pinned-original-source.tex, source lines 587--609. Cited by the retained statement or proof but not part of the collected spans; supplied only as a navigation pointer. The argument does not depend on its content.\par\medskip
\ifdefined\hypertarget\hypertarget{sourcepointer-3}{}\fi
\noindent\textit{Source reference.} Original theorem 01GY: sources/stacks-hypercovering/pinned-original-source.tex, source lines 1172--1196. Cited by the retained statement or proof but not part of the collected spans; supplied only as a navigation pointer. The argument does not depend on its content.\par\medskip
\ifdefined\hypertarget\hypertarget{sourcepointer-4}{}\fi
\noindent\textit{Source reference.} Original theorem 01GJ: sources/stacks-hypercovering/pinned-original-source.tex, source lines 1632--1642. Cited by the retained statement or proof but not part of the collected spans; supplied only as a navigation pointer. The argument does not depend on its content.\par\medskip
\ifdefined\hypertarget\hypertarget{sourcepointer-5}{}\fi
\noindent\textit{Source reference.} Original theorem 01GV: sources/stacks-hypercovering/pinned-original-source.tex, source lines 1062--1068. Cited by the retained statement or proof but not part of the collected spans; supplied only as a navigation pointer. The argument does not depend on its content.\par\medskip
\ifdefined\hypertarget\hypertarget{sourcepointer-6}{}\fi
\noindent\textit{Source reference.} Original theorem 01FW: sources/stacks-hypercovering/pinned-original-source.tex, source lines 2309--2309. Cited by the retained statement or proof but not part of the collected spans; supplied only as a navigation pointer. The argument does not depend on its content.\par\medskip
\begin{theorem}
\label{theorem-cohomology-hypercoverings}
Let $\mathcal{C}$ be a site with fibre products.
Let $X$ be an object of $\mathcal{C}$. Let $i \geq 0$.
The functors
\begin{eqnarray*}
\textit{Ab}(\mathcal{C}) & \longrightarrow & \textit{Ab} \\
\mathcal{F} & \longmapsto & H^i(X, \mathcal{F}) \\
\mathcal{F} & \longmapsto & \check{H}^i_{\text{HC}}(X, \mathcal{F})
\end{eqnarray*}
are canonically isomorphic.
\end{theorem}

\begin{proof}[Proof using spectral sequences.]
Suppose that $\xi \in H^p(X, \mathcal{F})$ for some $p \geq 0$.
Let us show that $\xi$ is in the image of the map
$\check{H}^p(X, \mathcal{F}) \to H^p(X, \mathcal{F})$ of
Lemma \ref{lemma-cech-spectral-sequence}
for some hypercovering $K$ of $X$.

\medskip\noindent
This is true if $p = 0$ by Lemma \ref{lemma-h0-cech}.
If $p = 1$, choose a {\v C}ech hypercovering $K$ of $X$ as in
Example \ref{example-cech} starting with a covering
$K_0 = \{U_i \to X\}$ in the site $\mathcal{C}$ such that
$\xi|_{U_i} = 0$, see
Cohomology on Sites,
Lemma \ref{sites-cohomology-lemma-kill-cohomology-class-on-covering}.
It follows immediately from the spectral sequence
in Lemma \ref{lemma-cech-spectral-sequence} that $\xi$ comes
from an element of $\check{H}^1(K, \mathcal{F})$ in this case.
In general, choose any hypercovering $K$ of $X$ such
that $\xi$ maps to zero in $\underline{H}^p(\mathcal{F})(K_0)$
(using Example \ref{example-cech} and
Cohomology on Sites,
Lemma \ref{sites-cohomology-lemma-kill-cohomology-class-on-covering}
again).
By the spectral sequence of Lemma \ref{lemma-cech-spectral-sequence}
the obstruction for $\xi$ to come from an element of
$\check{H}^p(K, \mathcal{F})$ is a sequence of elements
$\xi_1, \ldots, \xi_{p - 1}$ with
$\xi_q \in \check{H}^{p - q}(K, \underline{H}^q(\mathcal{F}))$
(more precisely the images of the $\xi_q$ in certain subquotients
of these groups).

\medskip\noindent
We can inductively replace the hypercovering $K$ by refinements
such that the obstructions $\xi_1, \ldots, \xi_{p - 1}$ restrict to zero
(and not just the images
in the subquotients -- so no subtlety here). Indeed, suppose we have
already managed to reach the situation where
$\xi_{q + 1}, \ldots, \xi_{p - 1}$ are zero.
Note that $\xi_q \in \check{H}^{p - q}(K, \underline{H}^q(\mathcal{F}))$
is the class of some element
$$
\tilde \xi_q \in
\underline{H}^q(\mathcal{F})(K_{p - q}) =
\prod H^q(U_i, \mathcal{F})
$$
if $K_{p - q} = \{U_i \to X\}_{i \in I}$. Let $\xi_{q, i}$
be the component of $\tilde \xi_q$ in $H^q(U_i, \mathcal{F})$.
As $q \geq 1$ we can use
Cohomology on Sites,
Lemma \ref{sites-cohomology-lemma-kill-cohomology-class-on-covering}
yet again to choose coverings $\{U_{i, j} \to U_i\}$
of the site such that each restriction $\xi_{q, i}|_{U_{i, j}} = 0$.
Consider the object $Z = \{U_{i, j} \to X\}$ of the category
$\text{SR}(\mathcal{C}, X)$ and its obvious morphism
$u : Z \to K_{p - q}$. It is clear that $u$ is a covering, see
Definition \ref{definition-covering-SR}. By
Lemma \ref{lemma-covering} there
exists a morphism $L \to K$ of hypercoverings of $X$ such that
$L_{p - q} \to K_{p - q}$ factors through $u$. Then clearly the
image of $\xi_q$ in $\underline{H}^q(\mathcal{F})(L_{p - q})$.
is zero. Since the spectral sequence of
Lemma \ref{lemma-cech-spectral-sequence}
is functorial this means that after replacing $K$ by $L$ we reach the
situation where $\xi_q, \ldots, \xi_{p - 1}$ are all zero.
Continuing like this we end up with a hypercovering where they are all
zero and hence $\xi$ is in the image of the map
$\check{H}^p(X, \mathcal{F}) \to H^p(X, \mathcal{F})$.

\medskip\noindent
Suppose that $K$ is a hypercovering of $X$, that
$\xi \in \check{H}^p(K, \mathcal{F})$ and that the image of
$\xi$ under the map
$\check{H}^p(X, \mathcal{F}) \to H^p(X, \mathcal{F})$ of
Lemma \ref{lemma-cech-spectral-sequence}
is zero. To finish the proof of the theorem we have to show that
there exists a morphism of hypercoverings $L \to K$ such that
$\xi$ restricts to zero in $\check{H}^p(L, \mathcal{F})$.
By the spectral sequence of Lemma \ref{lemma-cech-spectral-sequence}
the vanishing of the image of $\xi$ in $H^p(X, \mathcal{F})$
means that there exist elements $\xi_1, \ldots, \xi_{p - 2}$
with $\xi_q \in \check{H}^{p - 1 - q}(K, \underline{H}^q(\mathcal{F}))$
(more precisely the images of these in certain subquotients)
such that the images $d_{q + 1}^{p - 1 - q, q}\xi_q$ (in the spectral
sequence) add up to $\xi$. Hence by exactly the same mechanism as above
we can find a morphism of hypercoverings $L \to K$ such that
the restrictions of the elements $\xi_q$, $q = 1, \ldots, p - 2$
in $\check{H}^{p - 1 - q}(L, \underline{H}^q(\mathcal{F}))$ are zero.
Then it follows that $\xi$ is zero since the morphism $L \to K$
induces a morphism of spectral sequences according to
Lemma \ref{lemma-cech-spectral-sequence}.
\end{proof}

\end{document}
